\section{Gemini 4とAntigravityの利用議論}
Gemini 4 CarbonとHydrogenの使い分けや、Antigravityによるオーケストレーションについて利用者の投稿が続いた。HydrogenのSearchへの展開を肯定的に評価する声があり、「無料で全ユーザーに提供される」との言及も観測された。ただし、提供範囲の正確な条件は今回の収集では公式資料により確認されていない。開発者向けUIや機能への改善要望も見られる。元の公式発表は観測窓外であり、本号は窓内の利用反応を扱う。
\dailyxsource{利用者投稿}{https://x.com/MarceloLima/status/2122089162065285376}

\section{Claude Fable 5.2の価格・性能比較}
AnthropicのClaude Fable 5.2とMythos 5.2がOpenAIのモデルより安価だとする投稿が共有された。価格差を20～80\%とし、Fermatの最終定理の形式化など特定タスクでの高速化も主張している。しかし、Grokはこの情報をUnverifiedと判定しており、公式価格表やベンチマークの一次資料は確認されていない。単発投稿に基づく比較として扱い、製品の確定仕様とは区別する必要がある。
\dailyxsource{価格比較の投稿}{https://x.com/brainify78/status/2122088969618018698}

\section{NVIDIAとJensen Huangを巡る憶測}
Jensen HuangのNVIDIA復帰可能性について、予測市場の確率を引用した投稿や過去のオープンモデル支持発言に言及する投稿が見られた。復帰がAIのオープンソース動向に影響するとの見方も示されているが、公式声明は確認されていない。経営体制に関する憶測と、市場参加者の反応を記録したものであり、復帰予定や経営上の決定が確定したことを意味しない。
\dailyxsource{関連する憶測投稿}{https://x.com/TheNavroopSingh/status/2122051455997886662}

\section{ローカルAI評価ツールRealmark}
Matt MastracciがローカルAI向けベンチマークRealmarkの公開を投稿した。トークン生成速度（tok/s）だけでは利用体験を十分に表せないとして、タスク完了を重視するRIPALスコアを提案している。GLM 5.3 flashの実行レシピで品質が急変する例も紹介された。モデルや量子化設定の比較では速度と成果品質を併せて測る必要があるという問題提起であり、評価方法の詳細や再現性は今後の確認点となる。
\dailyxsource{開発者の発表投稿}{https://x.com/mmastrac/status/2122089331372544164}

\section{ワンショット動画生成の個人デモ}
ポストプロセスを使わないネイティブなワンショット動画生成をうたうデモ投稿が観測された。投稿者は、空間的な一貫性や複数の物理的相互作用を外部VFXなしで表現したと説明している。ただし使用モデル、制作条件、再現性は確認されておらず、Grokの判定はUnverifiedである。個人による技術デモとして紹介し、一般的な生成性能が実証されたものとは扱わない。
\dailyxsource{動画付きデモ投稿}{https://x.com/ilovevegas49341/status/2122089271293337842}

\section{オープンソースAIツールの再紹介}
Ollama、Open WebUI、CrewAIなどの既存ツールをまとめた投稿が共有された。ローカルモデルの実行、Web UIによる利用、複数エージェントの構築など、用途ごとの選択肢を有料サービスの代替として紹介する内容である。今回の観測は新しい機能公開ではなく、既存ツールの再共有に当たる。個々の機能や保守状況を新たに検証したものではない。
\dailyxsource{ツール一覧の投稿}{https://x.com/LucianoInsulaAi/status/2122086979592450543}

\section{Cursorのセキュリティ関連情報}
Cursorのセキュリティインシデントに関する更新に触れる投稿が観測された。ただしGrokが収集できたのは二次的な言及であり、Cursorの公式更新内容や対象範囲は確認されていない。具体的な脆弱性、被害、修正状況を断定する根拠は今回の入力にはない。開発支援AIの安全な運用に関わる話題として記録し、詳細は未確認とする。
\dailyxsource{二次的な関連投稿}{https://x.com/AI_tokusatsu/status/2122089232491835688}

\section{エージェントのループと課金管理}
AIエージェントが無限ループに陥り、高額な利用料金につながる事例があるとする注意喚起が投稿された。Grokの判定はUnverifiedで、具体的な請求事例や発生条件は検証されていない。一方で、長時間自律実行するエージェントのコスト管理という論点を示している。投稿は利用上のリスクに対する主張として紹介し、個別事例の事実認定は行わない。
\dailyxsource{コスト管理に関する投稿}{https://x.com/CodetocashAi/status/2122087122815312022}
